\section{Conclusions}

This paper presented SRPA-YOLO, a compact detector that formulates prediction-head capacity as task-specific representation allocation. A 64-channel shared representation supports classification and DFL localization, whereas a zero-output-initialized 16-channel classification-private representation contributes only an additive class residual. This construction increases the feature capacity available to classification while making the box output directly independent of the private representation. Exact structural re-parameterization converts the training-time decomposition into one 80-channel convolution and one merged output projection per scale, yielding a static single-path deployment graph.

On the deduplicated DUT Anti-UAV test set, the selected checkpoint achieves 95.850\% Precision, 92.201\% Recall, 95.185\% AP@0.50, 67.071\% \mapall{}, and 77.329\% AP@0.75 with 2.164 M deployed parameters. Compared with YOLOv8n, the corresponding improvements are 3.088, 8.242, 5.057, 6.895, and 9.750 percentage points, together with 28.0\% fewer parameters. The additional gains in Recall, AP@0.50, \mapall{}, and AP@0.75 over YOLOv8n-P2 show that head-level task allocation complements high-resolution prediction. The gains across the evaluated IoU thresholds and the low variation in the matched-seed experiments are consistent with stable localization and classification refinement under the reported protocol.

Two paired deployment measurements yield 9.7435 and 9.8574 ms for SRPA-YOLO and 9.8859 and 9.9893 ms for the paired YOLOv8n control under the specified RTX 5060, FP16, batch-1, forward-plus-NMS protocol. Together with the 28.0\% parameter reduction, these measurements indicate a favorable accuracy--compactness operating point on the evaluated GPU software stack. Future work should examine multi-class aerial scenes, embedded ONNX/TensorRT deployment, and scale-conditioned private capacity under additional target-size and imaging conditions.
